\section{Trabalhos Relacionados}
\label{sec_relacionados}

Esta seção apresenta os principais trabalhos e tecnologias relacionados ao serviço de metadados proposto neste trabalho. São discutidos o mecanismo de inicialização de instâncias cloud-init, os serviços de metadados das principais nuvens públicas e privadas, e as limitações atuais do Incus em relação a esse ecossistema.

\subsection{Cloud-init e Inicialização de Instâncias}

O cloud-init é o padrão da indústria para a inicialização automatizada de instâncias em ambientes de nuvem~\cite{cloudinit2024}. Desenvolvido originalmente pela Canonical e amplamente adotado pelos principais provedores de infraestrutura, o cloud-init é executado durante a inicialização do sistema operacional e é responsável por configurar aspectos fundamentais da instância, tais como nome de host, usuários e grupos, chaves SSH, pacotes de software, scripts de inicialização e parâmetros de rede.

O processo de inicialização do cloud-init é organizado em quatro estágios sequenciais, cada um com responsabilidades bem definidas~\cite{cloudinit2024}: (1)~\texttt{init-local}, executado antes da configuração de rede, responsável por identificar o \textit{datasource} a partir de fontes locais como disco, DMI e linha de comando do \textit{kernel}; (2)~\texttt{init} (rede), executado após a rede estar ativa, responsável por buscar os dados de configuração do \textit{datasource} identificado; (3)~\texttt{modules-config}, que aplica os módulos de configuração como criação de usuários e instalação de pacotes; e (4)~\texttt{modules-final}, que executa \textit{scripts} finais e mensagens de conclusão. Essa sequência é relevante porque impõe requisitos temporais à disponibilidade do serviço de metadados: os dados devem estar acessíveis antes do estágio \texttt{init}, momento em que o cloud-init realiza a busca efetiva dos metadados.

O funcionamento do cloud-init é baseado no conceito de \textit{datasources}: fontes de dados a partir das quais a ferramenta obtém as informações necessárias para configurar a instância. Cada datasource é composto por quatro tipos de dados distintos: (i) \textit{metadata}, que contém informações de identidade da instância como identificador único, nome de host e dados de rede; (ii) \textit{user-data}, que permite ao usuário final fornecer configurações personalizadas, como \textit{scripts} shell ou arquivos de configuração no formato cloud-config; (iii) \textit{vendor-data}, fornecido pelo provedor de infraestrutura para configurações específicas da plataforma; e (iv) \textit{network-config}, que define a configuração de rede da instância em formatos padronizados~\cite{cloudinit2024}. Dentre esses dados, o campo \texttt{instance-id} é particularmente crítico: o cloud-init utiliza-o para determinar se a inicialização corrente é a primeira (\textit{first boot}), acionando a configuração completa, ou uma inicialização subsequente, na qual módulos já aplicados são ignorados.

A capacidade do cloud-init de abstrair as particularidades de cada plataforma por meio da camada de datasources é o que o torna uma solução portável e amplamente adotada. No entanto, para que essa abstração funcione corretamente, é necessário que a infraestrutura subjacente disponibilize um endpoint de metadados que o cloud-init possa consultar durante a inicialização.

\subsection{Datasource NoCloud}

O datasource \textit{NoCloud} é o mais flexível dentre os oferecidos pelo cloud-init e é diretamente relevante para ambientes de infraestrutura privada~\cite{cloudinit_nocloud}. Ele permite obter dados de configuração a partir de uma URL HTTP base (\texttt{seedfrom}), consultando os caminhos \texttt{/meta-data}, \texttt{/user-data}, \texttt{/vendor-data} e \texttt{/network-config} relativos a essa URL. Essa abordagem não exige que o serviço implemente a hierarquia de caminhos versionados do EC2 (como \texttt{/latest/meta-data/}), bastando que os quatro endpoints estejam disponíveis na raiz configurada. O \textit{NoCloud} suporta os protocolos HTTP, HTTPS, FTP e FTPS, e requer apenas que o arquivo \texttt{meta-data} contenha, no mínimo, os campos \texttt{instance-id} e \texttt{local-hostname}~\cite{cloudinit_nocloud}. Essa simplicidade torna o \textit{NoCloud} especialmente adequado para serviços de metadados independentes de plataforma, como o proposto neste trabalho.

\subsection{AWS Instance Metadata Service (IMDS)}

A Amazon Web Services (AWS) disponibiliza o \textit{Instance Metadata Service} (IMDS), acessível pelas instâncias EC2 por meio do endereço IP de link-local \texttt{169.254.169.254}~\cite{aws_imds}. Esse serviço expõe uma API HTTP que permite às instâncias consultar informações sobre si mesmas sem necessidade de credenciais adicionais.

O IMDS oferece duas versões com características distintas de segurança. A primeira versão, IMDSv1, opera com requisições GET simples sem autenticação, o que pode expor a instância a ataques de falsificação de requisição do lado do servidor (\textit{Server-Side Request Forgery} — SSRF)~\cite{mitre_t1552_005}. A segunda versão, IMDSv2, introduzida em 2019, adota um modelo baseado em sessões: a instância deve primeiro obter um token de acesso por meio de uma requisição PUT com cabeçalho \texttt{X-aws-ec2-metadata-token-ttl-seconds}, e em seguida utilizar esse token nas requisições subsequentes via o cabeçalho \texttt{X-aws-ec2-metadata-token}~\cite{aws_imdsv2}. Essa abordagem mitiga significativamente os riscos associados a ataques SSRF, uma vez que a maioria das vulnerabilidades SSRF permite apenas requisições GET e não possibilita a inclusão de cabeçalhos personalizados. Além disso, o IMDSv2 aplica um limite de saltos (\textit{hop limit}) de 1 nas respostas PUT e rejeita requisições contendo o cabeçalho \texttt{X-Forwarded-For}, bloqueando ataques baseados em proxies reversos~\cite{aws_imdsv2_blog}.

Os dados servidos pelo IMDS incluem, entre outros: identificador único da instância (\texttt{instance-id}), nome de host, endereços MAC das interfaces de rede, metadados de interfaces de rede, chaves SSH públicas associadas à instância e o \textit{user-data} fornecido no momento do lançamento. A API é organizada em um namespace hierárquico acessível via caminhos como \texttt{/latest/meta-data/} e \texttt{/latest/user-data}.

\subsection{GCP Metadata Server}

O Google Cloud Platform (GCP) disponibiliza o \textit{Metadata Server}, acessível pelo nome de host \texttt{metadata.google.internal} (resolvido para \texttt{169.254.169.254}) e também diretamente pelo endereço IP \texttt{metadata.google.internal}~\cite{gcp_metadata}. O serviço é organizado em dois escopos principais: metadados de projeto, que contêm informações compartilhadas entre todas as instâncias de um projeto, e metadados de instância, específicos de cada máquina virtual.

Um aspecto de segurança relevante do Metadata Server do GCP é a exigência do cabeçalho HTTP \texttt{Metadata-Flavor: Google} em todas as requisições. Requisições que não incluam esse cabeçalho são rejeitadas pelo serviço, o que serve como mecanismo de proteção contra ataques SSRF oriundos de contextos externos à instância~\cite{gcp_metadata}. Mecanismo semelhante é adotado pelo Azure Instance Metadata Service, que exige o cabeçalho \texttt{Metadata: true}~\cite{azure_imds}. Embora esses cabeçalhos obrigatórios mitiguem ataques SSRF simples — que tipicamente não permitem a inclusão de cabeçalhos personalizados — oferecem proteção inferior ao modelo de sessões do IMDSv2 da AWS~\cite{aws_imdsv2_blog}. Além da proteção contra SSRF, o serviço do GCP suporta notificações de alterações de metadados via mecanismo de \textit{long polling}, permitindo que aplicações reajam dinamicamente a mudanças de configuração.

\subsection{OpenStack Metadata Service}

O OpenStack, principal plataforma de nuvem privada de código aberto, implementa um serviço de metadados compatível com a API EC2 da AWS, também disponível no endereço \texttt{169.254.169.254}~\cite{openstack_metadata}. O serviço oferece dois endpoints distintos: o endpoint compatível com EC2 (\texttt{/latest/meta-data/}) e o endpoint nativo do OpenStack (\texttt{/openstack/}), que serve dados em formato JSON estruturado conforme o modelo de dados da plataforma.

Uma característica fundamental do serviço de metadados do OpenStack é seu forte acoplamento com os demais componentes da plataforma. O roteamento de requisições de metadados é realizado pelo Neutron, o componente de rede do OpenStack, que intercepta requisições destinadas ao endereço \texttt{169.254.169.254} e as encaminha ao serviço correto com base na interface de rede da instância solicitante. O provisionamento dos dados de metadados, por sua vez, é gerenciado pelo Nova, o componente de computação. Esse acoplamento implica que o serviço de metadados do OpenStack não pode ser executado de forma autônoma fora do ecossistema OpenStack~\cite{openstack_metadata}, o que limita sua aplicabilidade em outros contextos de virtualização. Em contextos acadêmicos, a dependência do ecossistema completo é aceita quando a escala justifica: Bollig et al.~\cite{bollig2018leveraging} descrevem a plataforma Stratus, uma nuvem privada baseada em OpenStack na Universidade de Minnesota, na qual o cloud-init consulta o serviço de metadados do Nova durante a inicialização para configurar endereços IP, chaves SSH e, por meio do \textit{vendor-data}, aplicar atualizações de segurança obrigatórias em todas as máquinas virtuais. Esse exemplo ilustra como o padrão de metadata service é fundamental para operações de nuvem privada, mas também evidencia que sua adoção requer a implantação de uma plataforma completa como o OpenStack.

\subsection{Incus/LXD e a Lacuna de Metadados}

O Incus é um gerenciador moderno de contêineres de sistema e máquinas virtuais, surgido em 2023 como um \textit{fork} da comunidade do projeto LXD, originalmente desenvolvido e mantido pela Canonical~\cite{stgraber2023, incus2024}. O projeto foi iniciado por Stéphane Graber, criador original do LXD, após a Canonical encerrar a participação comunitária no desenvolvimento do projeto~\cite{stgraber2023}. O Incus herda a arquitetura e as funcionalidades do LXD, com o objetivo de manter o desenvolvimento aberto e orientado pela comunidade~\cite{lxd2024}.

O Incus possui suporte nativo para configurações do cloud-init por meio de chaves de configuração específicas: \texttt{cloud-init.user-data} para dados de usuário, \texttt{cloud-init.vendor-data} para dados do fornecedor e \texttt{cloud-init.network-config} para configuração de rede~\cite{incus2024}. Essas configurações podem ser definidas diretamente na instância ou propagadas via perfis de configuração. No entanto, o Incus não disponibiliza nenhum endpoint HTTP de metadados que as instâncias possam consultar durante a inicialização.

A entrega das configurações do cloud-init no Incus ocorre por meio do \textit{datasource LXD}, que se comunica com o hipervisor por um \textit{socket} Unix em \texttt{/dev/incus/sock} (ou \texttt{/dev/lxd/sock} no LXD original)~\cite{cloudinit_lxd_ds}. Esse \textit{socket} é provido pelo \textit{Incus agent}, um processo leve executado dentro das instâncias. Alternativamente, as configurações podem ser injetadas diretamente no sistema de arquivos da instância via \textit{config drive}~\cite{incus2024}. Embora esses mecanismos sejam funcionais para casos de uso básicos, eles apresentam uma limitação temporal crítica: o \textit{Incus agent} é iniciado tardiamente na sequência de \textit{boot} da máquina virtual, de modo que os estágios iniciais do cloud-init (\texttt{ds-identify} e \texttt{init-local}) frequentemente são executados antes que o \textit{socket} esteja disponível. Essa condição de corrida resulta em falhas na detecção do datasource e, consequentemente, na configuração automática da instância.

Trabalhos como o de Sousa~\cite{sousa2018computacao}, que implementa uma infraestrutura de nuvem privada na USP utilizando LXD com OpenNebula para instanciação de contêineres sob demanda, evidenciam que a plataforma é viável para uso acadêmico e de pesquisa, mas a inicialização das instâncias depende de \textit{scripts} definidos manualmente nos \textit{templates} da plataforma de orquestração, sem um serviço de metadados padronizado.

Além dessa limitação temporal, os mecanismos atuais do Incus não replicam o fluxo completo do datasource do cloud-init que opera via requisições HTTP ao endereço \texttt{169.254.169.254}. Isso impede o uso do datasource \textit{NoCloud} HTTP~\cite{cloudinit_nocloud} ou do datasource \textit{EC2}~\cite{cloudinit_ec2_ds}, que são os mais amplamente suportados e testados no ecossistema cloud-init, e limita a compatibilidade com imagens e ferramentas de automação que esperam um serviço de metadados acessível via rede.

\subsection{Algoritmo de Consenso RAFT}

O RAFT é um algoritmo de consenso distribuído projetado com ênfase na compreensibilidade, proposto por Ongaro e Ousterhout~\cite{ongaro2014raft}. Ao contrário do Paxos, que é reconhecidamente difícil de compreender e implementar corretamente, o RAFT decompõe o problema do consenso em subproblemas relativamente independentes: eleição de líder, replicação de log e segurança. Essa decomposição facilita tanto o entendimento do algoritmo quanto a construção de implementações corretas.

No RAFT, um cluster é composto por um nó líder e um ou mais seguidores. O líder é responsável por receber requisições de escrita, replicá-las para os seguidores via entradas de log, e confirmar a operação assim que uma maioria de nós tiver persistido a entrada. Essa garantia de quórum assegura que o sistema tolera falhas de até $\lfloor(n-1)/2\rfloor$ nós em um cluster de $n$ membros, mantendo disponibilidade e consistência.

A biblioteca \texttt{hashicorp/raft}~\cite{hashicorp_raft} é uma implementação em Go amplamente utilizada do protocolo RAFT, desenvolvida e mantida pela HashiCorp. Ela fornece uma API estável para integração do consenso distribuído em aplicações, sendo empregada em projetos de destaque como o Consul e o Vault. Neste trabalho, a biblioteca é utilizada para prover alta disponibilidade ao serviço de metadados, garantindo que o estado do cluster permaneça consistente mesmo diante de falhas de nós individuais.

\subsection{Segurança em Serviços de Metadados}

Serviços de metadados representam uma superfície de ataque relevante em ambientes de nuvem, uma vez que expõem informações sensíveis — como credenciais temporárias, chaves SSH e configurações de rede — por meio de um endpoint HTTP acessível sem autenticação convencional. O vetor de ataque mais significativo é o \textit{Server-Side Request Forgery} (SSRF): caso uma aplicação web executada dentro da instância possua uma vulnerabilidade que permita a um atacante induzir requisições HTTP arbitrárias, o atacante pode direcionar essas requisições ao endereço \texttt{169.254.169.254} e obter os metadados da instância~\cite{mitre_t1552_005}. Esse vetor é catalogado pela MITRE ATT\&CK como a técnica T1552.005 (\textit{Unsecured Credentials: Cloud Instance Metadata API}).

Um caso emblemático que ilustra a gravidade dessa classe de ataque é a violação de dados da Capital One, ocorrida em março de 2019 e analisada em profundidade por Novaes Neto et al.~\cite{novaes2020capital}. Nesse incidente, um atacante explorou uma vulnerabilidade SSRF em um \textit{Web Application Firewall} (WAF) mal configurado para induzir requisições ao IMDSv1 da AWS no endereço \texttt{169.254.169.254}, obtendo credenciais temporárias IAM associadas à instância. Com essas credenciais, o atacante acessou buckets S3 contendo dados de 106~milhões de clientes, resultando em prejuízos superiores a 150~milhões de dólares. A cadeia de ataque --- SSRF, acesso ao IMDS sem autenticação, exfiltração de credenciais --- demonstrou que a simplicidade do IMDSv1, baseado em requisições GET sem qualquer verificação, representava um risco sistêmico. Esse incidente foi um dos catalisadores para a introdução do IMDSv2 pela AWS.

Os provedores de nuvem adotam diferentes estratégias de mitigação. A AWS introduziu o IMDSv2 com autenticação baseada em sessões~\cite{aws_imdsv2_blog}, o GCP exige o cabeçalho \texttt{Metadata-Flavor: Google}~\cite{gcp_metadata}, e o Azure requer o cabeçalho \texttt{Metadata: true}~\cite{azure_imds}. Em ambientes de infraestrutura privada, como os atendidos pelo serviço proposto neste trabalho, o risco de SSRF é tipicamente menor, pois o serviço de metadados não está exposto à internet e o perímetro de rede é controlado pelo administrador. Ainda assim, a documentação do risco e a adoção de boas práticas — como a restrição de acesso por endereço IP de origem e a possibilidade futura de autenticação por token — são medidas relevantes para ambientes de produção.

\subsection{Comparação entre as Abordagens}

A Tabela~\ref{tab:comparacao} sintetiza as principais características dos serviços de metadados discutidos nesta seção em comparação com a proposta deste trabalho. Os critérios avaliados são: operação autônoma (sem dependência de plataforma), compatibilidade com o cloud-init via endpoint HTTP, suporte a contêineres de sistema, suporte a máquinas virtuais, alta disponibilidade e disponibilidade como software de código aberto.

\begin{table}[ht!]
\centering
\small
\setlength{\tabcolsep}{5pt}
\renewcommand{\arraystretch}{1.3}
\caption{Comparação entre serviços de metadados para cloud-init.}
\vspace{0.5em}
\label{tab:comparacao}
\def\cellZ{\cellcolor{black!10}\xmark}
\def\cellO{\cmark}
\begin{tabular}{>{\raggedright\arraybackslash}p{3.2cm} >{\centering\arraybackslash}p{1.8cm} >{\centering\arraybackslash}p{2.0cm} >{\centering\arraybackslash}p{1.8cm} >{\centering\arraybackslash}p{1.6cm} >{\centering\arraybackslash}p{1.5cm} >{\centering\arraybackslash}p{1.8cm}}
    \toprule
    \textbf{Solução} & \textbf{Autônomo} & \textbf{Compat. cloud-init} & \textbf{Suporta Contêineres} & \textbf{Suporta VMs} & \textbf{Alta Disp.} & \textbf{Código Aberto} \\
    \midrule
    AWS IMDS            & \cellZ & \cellO & \cellZ & \cellO & \cellO & \cellZ \\
    GCP Metadata        & \cellZ & \cellO & \cellZ & \cellO & \cellO & \cellZ \\
    OpenStack Metadata  & \cellZ & \cellO & \cellZ & \cellO & \cellO & \cellO \\
    LXD/Incus (atual)   & \cellO & \cellZ & \cellO & \cellO & \cellZ & \cellO \\
    \textbf{Este trabalho} & \cellO & \cellO & \cellO & \cellO & \cellO & \cellO \\
    \bottomrule
\end{tabular}
\end{table}

\subsection{Considerações Finais}

A análise dos trabalhos relacionados evidencia uma lacuna tecnológica relevante: não existe atualmente um serviço de metadados leve e autônomo para o Incus que implemente o protocolo HTTP de metadados esperado pelo cloud-init, habilitando o fluxo completo de datasource para contêineres e máquinas virtuais gerenciados por essa plataforma. Os serviços existentes nas nuvens públicas (AWS IMDS e GCP Metadata Server) são proprietários e estão acoplados às respectivas infraestruturas. O serviço de metadados do OpenStack, apesar de ser código aberto e compatível com cloud-init, depende integralmente do ecossistema OpenStack para funcionar, inviabilizando seu uso isolado. O próprio Incus, por sua vez, oferece mecanismos de entrega de configuração baseados no datasource LXD~\cite{cloudinit_lxd_ds} que, além de não cobrirem o fluxo completo de consulta HTTP durante a inicialização, sofrem de uma condição de corrida temporal durante a inicialização. A necessidade de ferramentas para inicialização automatizada de instâncias em nuvens de infraestrutura é documentada na literatura desde 2011, quando Bresnahan et al.~\cite{bresnahan2011managing} propuseram o cloudinit.d, uma ferramenta para lançamento coordenado de conjuntos de máquinas virtuais interdependentes em nuvens IaaS, evidenciando já naquele período a dificuldade de configurar instâncias de forma repetível e automatizada sem o suporte de serviços de metadados. A adoção crescente de práticas de Infraestrutura como Código (\textit{Infrastructure as Code} --- IaC)~\cite{guerriero2019iac, artac2017iac} reforça a importância de serviços de metadados padronizados para o provisionamento automatizado de instâncias. Nesse contexto, Teppan et al.~\cite{teppan2022survey} observam que, dentre as categorias de ferramentas \textit{cloud-native} catalogadas pela CNCF, a categoria de provisionamento é a que menos dispõe de projetos maduros para ambientes \textit{on-premise}, o que corrobora a lacuna identificada neste trabalho para o ecossistema Incus.

Este trabalho propõe um serviço de metadados compatível com cloud-init, projetado especificamente para ambientes Incus, que opera de forma autônoma, suporta tanto contêineres quanto máquinas virtuais, oferece alta disponibilidade via o algoritmo de consenso RAFT~\cite{ongaro2014raft} implementado pela biblioteca \texttt{hashicorp/raft}~\cite{hashicorp_raft}, e é disponibilizado como software de código aberto. Dessa forma, preenche a lacuna identificada e viabiliza o uso do ecossistema completo do cloud-init em infraestruturas baseadas em Incus.
